\documentclass[../../../include/open-logic-section]{subfiles}

\begin{document}

\olfileid{sth}{cardinals}{hp}
\olsection{পরিশিষ্ট: হিউমের নীতি}

\olref[cp]{sec}-এ আমরা কান্টরের নীতি বলেছিলাম। সেটি ছিল:
\begin{align*}
	\card{A} = \card{B} & \text{ iff } A \approx B.
\intertext{এর সঙ্গে বর্তমানে \emph{হিউমের নীতি} নামে
পরিচিত নীতিটির খুব মিল আছে। সেটি বলে:}
	\fregenum{x} {F(x)} = \fregenum{x}{G(x)} & \text{ iff } F \sim G
\end{align*}
এখানে `$F\sim G$' বলতে বোঝায়, $F$ পূরণকারী বস্তুর সংখ্যা ও $G$ পূরণকারী বস্তুর
সংখ্যা ঠিক সমান; অর্থাৎ $F$-গুলি এবং $G$-গুলির
মধ্যে একটি একৈক ও সমাপতিত মিল স্থাপন করা যায়।
অন্যভাবে:
\begin{align*}
	\exists R(&\forall v\forall y(Rvy \lif (Fv \land Gy)) \land {}\\
		&\forall v(Fv \lif \lexists![y][Rvy]) \land {}\\
		&\forall y(Gy \lif \lexists![v][Rvy]))
\end{align*}
কিন্তু হিউমের নীতি ও কান্টরের নীতির মধ্যে প্রকারগত
পার্থক্য আছে। কান্টরের নীতিতে ``$A$'' ও ``$B$''
প্রথম-ক্রমের পদ, যেগুলি \emph{সেট} নির্দেশ করে।
হিউমের নীতিতে ``$F$'', ``$G$'' ও ``$R$''
প্রথম-ক্রমের পদ \emph{নয়}; তারা \emph{বিধেয়ের
স্থানে} আছে। (সম্ভবত তারা \emph{ধর্ম} নির্দেশ করে।)
তাই হিউমের নীতিকে কথায় এভাবে বলা যায়:
$F$-গুলির সংখ্যা $G$-গুলির সংখ্যার সমান তখনই,
যখন দুই শ্রেণির বস্তুর মধ্যে একৈক ও সমাপতিত মিল
স্থাপন করা যায়। এখানে $F$-গুলি ও $G$-গুলির
সংখ্যাই তুলনা করছি। একে \emph{হিউমের নীতি} বলা হয়,
কারণ হিউম একসময় লিখেছিলেন:
\begin{quote}
  দুটি সংখ্যা এমনভাবে মেলানো হলে, যাতে একটির প্রত্যেকটি
  এককের বিপরীতে অন্যটির একটি একক থাকে, আমরা তাদের
  সমান বলি। \citep[প্রথম পুস্তক, তৃতীয় ভাগ, \S1]{Hume1740}
\end{quote}
\citet[\S63]{Frege1884} হিউমের এই কথাটি উদ্ধৃত
করে পরে, কার্যত, আমরা যাকে হিউমের নীতি বলেছি
তার একটি রূপরেখা দেন। এর ফলে সমকালীন গাণিতিক
যুক্তিবিদ্যায় নীতিটি গুরুত্ব পায়।

হিউমের নীতি ও মৌলিক সূত্র~V-এর কাঠামোগত মিল লক্ষ করুন।
\olref[story][blv]{sec}-এ সূত্রটি এভাবে লিখেছিলাম:
\[
	\fregeext{x}{F(x)} = \fregeext{x}{G(x)} \text{iff } \lforall[x][(F(x) \liff G(x))].
\]
এখানে তাদের তুলনা কিছুটা খুলে বলা দরকার।

হিউমের নীতি কিংবা মৌলিক সূত্র~V-এর মতো নীতিকে
দুইভাবে বোঝা যায়: \emph{প্রেডিকেটিভ} অথবা
\emph{ইমপ্রেডিকেটিভ} অর্থে
(মনে করুন \olref[story][predicative]{sec})।
মৌলিক সূত্র~V-এর ইমপ্রেডিকেটিভ পাঠে, প্রত্যেক $F$-এর
জন্য $\fregeext{x}{F(x)}$ বস্তুটি সেই পরিমাণায়নের
ক্ষেত্রেই পড়ে, যা দিয়ে সূত্র~V বলা হয়েছে।
একইভাবে হিউমের নীতির ইমপ্রেডিকেটিভ পাঠে,
প্রত্যেক $F$-এর জন্য $\fregenum{x}{F(x)}$
বস্তুটিও নীতিটি বলার সময় ব্যবহৃত পরিমাণায়নের
ক্ষেত্রেই পড়ে। বিপরীতে \emph{প্রেডিকেটিভ} পাঠে
$\fregeext{x}{F(x)}$ ও $\fregenum{x}{F(x)}$
বস্তুগুলি অন্য কোনো ক্ষেত্রের সদস্য হবে।

মৌলিক সূত্র~V ইমপ্রেডিকেটিভ অর্থে নিলে অনানুষ্ঠানিক ধর্মনির্দেশে
সেট-গঠনের মাধ্যমে অসংগতি দেখা দেয়
(বিশদ দেখুন \olref[story][blv]{sec})। অনানুষ্ঠানিক ধর্মনির্দেশে সেট-গঠনের
মতোই প্রেডিকেটিভ পাঠে একে সংগতি-সম্মত করা যায়;
কিন্তু তাতে আমরা যা চাই, তার সবই সম্ভবত মিলবে না।

অন্যদিকে হিউমের নীতি \emph{ইমপ্রেডিকেটিভ} অর্থেও
সংগতভাবে নেওয়া যায়, এবং সে অর্থে নীতিটি
বেশ শক্তিশালী।
% BN-SRC-807 TARGET-CORRECTION-BEGIN
\emph{উৎস-সংশোধন-নোট।} আনুষ্ঠানিক পরিধি এখানে
পরিষ্কার রাখা দরকার। পূর্ণ বিধেয়/সম্পর্ক-ক্ষেত্রের সঙ্গে
মৌলিক সূত্র~V-র বিরোধ; সম্পর্ক-ক্ষেত্র সীমিত করলে
শুধু নামের ইমপ্রেডিকেটিভত্ব থেকে বিরোধ অনুসৃত নয়।
হিউমের নীতির ক্ষেত্রে বাইরের সেটতত্ত্বে একটি স্পষ্ট মডেল
আছে: বস্তু-ক্ষেত্র $D=\omega+1$, সমস্ত বিধেয়
$\Pow D$ এবং সমস্ত দ্বিপদী সম্পর্ক $\Pow{D\times D}$।
সসীম $F$-এ তার স্বাভাবিক সদস্যসংখ্যা এবং অসীম $F$-এ
$\omega$ বস্তুটিকে সংখ্যা নিই। এগুলি সব $D$-সদস্য;
দুই অসীম উপসেট গণনীয়, আর দুই সসীম উপসেটে সংখ্যা
একই ঠিক তখনই যখন একৈক ও সমাপতিত চিত্রণ আছে। ফলে পুরো নীতি
সত্য। এটি বাইরের মডেল-নির্মাণের সাপেক্ষে সঙ্গততার
ব্যাখ্যা; বাইরের তত্ত্বের নিজস্ব সঙ্গততা প্রমাণ নয়।
% BN-SRC-807 TARGET-CORRECTION-END

উদাহরণ হিসেবে ``$x\neq x$'' বিধেয়টি ধরুন;
কোনো বস্তুই এটি পূরণ করে না। হিউমের নীতি থেকে
এখন $\#x(x\neq x)$ নামে একটি বস্তু পাই। একে
$0$ সংখ্যা ধরতে পারি। \emph{ইমপ্রেডিকেটিভ} পাঠে,
এবং \emph{শুধু} সেই পাঠেই, এই $0$ আমাদের মূল
পরিমাণায়নের ক্ষেত্রের মধ্যে পড়ে। তাই
``$x=0$'' বিধেয়ের উপরও হিউমের নীতি প্রয়োগ করে
$\#x(x=0)$ পাই; একে $1$ সংখ্যা ধরতে পারি।
হিউমের নীতি থেকে $0\neq 1$-ও পাওয়া যায়:
নিজের থেকে ভিন্ন বস্তু এবং $0$-র সমান বস্তুর মধ্যে
একৈক ও সমাপতিত মিল সম্ভব নয়। প্রথম ধরনের
বস্তু একটিও নেই, দ্বিতীয় ধরনের একটি আছে।
ইমপ্রেডিকেটিভ পাঠে এবার $1$-ও আমাদের মূল
পরিমাণায়নের ক্ষেত্রের সদস্য। তাই
``$(x=0\lor x=1)$'' বিধেয়ের জন্য হিউমের নীতি
প্রয়োগ করে $\#x(x=0\lor x=1)$ পাই;
একে $2$ সংখ্যা ধরতে পারি। একইভাবে
$0\neq 2$, $1\neq 2$ ইত্যাদিও দেখানো যায়।

সুতরাং ইমপ্রেডিকেটিভ অর্থে হিউমের নীতি থেকে
\emph{অসীমসংখ্যক বস্তুর} অস্তিত্ব মেলে। এই কারণে
\emph{নব্য-ফ্রেগীয় যুক্তিবাদীরা} পাটিগণিতের ভিত্তি
হিসেবে হিউমের নীতি গ্রহণে উৎসাহ পেয়েছেন।

তবে ফ্রেগে \emph{নিজে} হিউমের নীতিকে পাটিগণিতের
ভিত্তি করেননি। বরং একটি স্পষ্ট সংজ্ঞা থেকে নীতিটি
প্রমাণ করেছিলেন: $\fregenum{x}{F(x)}$-কে
$F\sim\Phi$ ধারণার বিস্তার হিসেবে সংজ্ঞায়িত করা হয়।
আধুনিক পরিভাষায় একে
$\fregenum{x}{F(x)}=\Setabs{G}{F\sim G}$
লেখার চেষ্টা করতে পারি; কিন্তু তাতে আবার অনানুষ্ঠানিক ধর্মনির্দেশে
সেট-গঠনের সমস্যায় ফিরে যাই।

\end{document}
